\documentclass{article} 
\input{../../lib.sty} 
 
 
\title{\texttt{fn make\_randomized\_response}} 
\author{\MichaelShoemateAuthor} 
\begin{document} 
 
\maketitle 
 
Proves soundness of \rustdoc{measurements/fn}{make\_randomized\_response} in \asOfCommit{mod.rs}{f5bb719}, 
a constructor taking a category set \texttt{categories} and probability \texttt{prob}. 
The mechanism returned by \texttt{make\_randomized\_response} takes in a data set \texttt{arg} (a single category), and... 
 
\begin{itemize} 
    \item ...if \texttt{arg} is in \texttt{categories}, 
    the mechanism truthfully returns the same value \texttt{arg} with probability \texttt{prob}, 
    otherwise it lies by selecting one of the other categories uniformly at random. 
    \item ...if \texttt{arg} is not in \texttt{categories},  
    it returns a category chosen uniformly at random. 
\end{itemize} 
 
\section{Hoare Triple} 
 
\subsection*{Precondition} 

\subsubsection*{Compiler-verified} 
\begin{itemize} 
    \item Variable \texttt{categories} must be a set with members of type \texttt{T} 
    \item Variable \texttt{prob} must be of type \texttt{f64} 
\end{itemize} 

\subsubsection*{Caller-verified} 
None 
 
\subsection*{Pseudocode} 
\lstinputlisting[language=Python,firstline=2,escapechar=|]{./pseudocode/make_randomized_response.py} 
 
\subsection*{Postcondition} 
 
\begin{theorem} 
\validMeasurement{\texttt{(categories, prob, T)}}{\\ \texttt{make\_randomized\_response}} 
\end{theorem} 
 
\begin{proof}  
 
\begin{tcolorbox} 
    The proof assumes that \texttt{sample\_uniform\_uint\_below} and 
    \texttt{sample\_bernoulli\_float} satisfy their postconditions. 
\end{tcolorbox} 
 
\texttt{sample\_uniform\_uint\_below} and \texttt{sample\_bernoulli\_float} can only fail due to lack of system entropy. 
This is usually related to the computer's physical environment and not the dataset.  
The rest of this proof is conditioned on the assumption that these functions do not raise an exception.  
 
Let $x$ and $x'$ be datasets that are \texttt{d\_in}-close with respect to \texttt{input\_metric}. 
Here, the metric is \texttt{DiscreteMetric} which enforces that $\din \geq 1$ if $x \ne x'$ and $\din = 0$ if $x = x'$.  
If $x = x'$, then the output distributions on $x$ and $x'$ are identical, and therefore the max-divergence is 0. 
 
Now consider the case where $x \ne x'$.  
For shorthand, we let $p$ represent \texttt{prob}, the probability of returning the input, 
and $t$ denote the number of categories. 
The implementation first converts the input set \texttt{categories} into a vector. 
Because the input data type is a set, each category appears exactly once, so this conversion only fixes an arbitrary order and does not change membership or multiplicity. 
 
$t$ must be at least two, by pseudocode line \ref{line:num_cats}, as any fewer would not be useful. 
$p$ is restricted to $[1/t, 1.0]$ by pseudocode line \ref{line:range}, as any less would not be useful. 
 
If $p = 1$, then pseudocode line \ref{line:map} sets the privacy constant to $\infty$. 
This matches the implementation. In this case, whenever $x$ and $x'$ are two distinct members of \texttt{categories}, 
the mechanism returns $x$ and $x'$ deterministically, so the max-divergence is $\infty$. 
Therefore the privacy bound holds when $p = 1$. For the remainder of the proof, assume $p < 1$. 
 
We now characterize the output distribution of the mechanism. 
If $x \in \texttt{categories}$, then pseudocode line \ref{line:fn} sets \texttt{index} 
to the unique position of $x$ in the vectorized category list. 
The call to \texttt{sample\_uniform\_uint\_below} samples uniformly from 
$\{0, \ldots, t - 2\}$. If the sampled index is greater than or equal to 
\texttt{index}, the implementation increments it by one, so the final index is 
uniform over the $t - 1$ positions corresponding exactly to 
\texttt{categories} $\setminus \{x\}$. Therefore \texttt{lie} is uniformly distributed over 
the categories other than $x$. The call to \texttt{sample\_bernoulli\_float(prob, False)} 
returns \texttt{true} with probability $p$, so the mechanism returns $x$ with 
probability $p$ and otherwise returns a uniformly sampled lie. 

If $x \not\in \texttt{categories}$, then \texttt{index} is \texttt{None}, so 
\texttt{sample\_uniform\_uint\_below} samples uniformly from all $t$ category 
indices and the shift is skipped. In this branch, \texttt{is\_member} is 
\texttt{false}, so the mechanism always returns \texttt{lie}. Therefore the 
output is uniformly distributed over \texttt{categories}. 
 
For any outcome $y \in \texttt{categories}$, the probability of observing $y$ is 
one of three values: 
 
\begin{enumerate} 
    \item If $x \in \texttt{categories}$ and $y = x$: 
    \[ 
        \Pr[M(x) = y \mid y = x] = p 
    \] 
 
    \item If $x \in \texttt{categories}$ and $y \ne x$: 
    \[ 
        \Pr[M(x) = y \mid y \ne x \wedge x \in \texttt{categories}] = \frac{1 - p}{t - 1} 
    \] 
 
    \item If $x \not\in \texttt{categories}$: 
    \[ 
        \Pr[M(x) = y \mid x \not\in \texttt{categories}] = \frac{1}{t} 
    \] 
\end{enumerate} 
 
\begin{tcolorbox} 
    \textbf{Claim.} The probability of case 3 is bounded by cases one and two: 
    \begin{equation} 
        \label{eq:bounded-case-3} 
        \frac{1 - p}{t - 1} \leq \frac{1}{t} \leq p 
    \end{equation} 

    The upper bound follows from pseudocode line \ref{line:range}, which gives 
    $p \ge 1/t$. Reusing \ref{line:range}, 
    \[ 
        \frac{1 - p}{t - 1} \leq \frac{1 - 1/t}{t - 1} = \frac{1}{t}. 
    \] 
    Therefore $1/t$ is also bounded below by case two 
    ($\frac{1 - p}{t - 1}$). 
\end{tcolorbox} 
 
By \eqref{eq:bounded-case-3}, every point probability lies between 
$\frac{1 - p}{t - 1}$ and $p$. Therefore, for any neighboring $x, x'$ and any 
$y \in \texttt{categories}$, 
\[ 
    \frac{\Pr[M(x) = y]}{\Pr[M(x') = y]} 
    \leq \frac{p}{(1 - p) / (t - 1)} 
    = \frac{p \cdot (t - 1)}{1 - p}. 
\] 
 
We now consider the max-divergence of the mechanism over all choices of neighboring datasets. 
 
\begin{align} 
    &\max_{x \sim x'} D_{\infty}(M(x), M(x')) \\ 
    =& \max_{x \sim x'} \max_{S \subseteq Supp(M(x))}\ln (\frac{\Pr[M(x) \in S]}{\Pr[M(x') \in S]}) \\ 
    \le& \max_{x \sim x'} \max_{y \in Supp(M(x))}\ln (\frac{\Pr[M(x) = y]}{\Pr[M(x') = y]}) &&\text{Lemma 3.3 } \cite{Kasiviswanathan_2014} \\ 
    \le& \ln \left(\frac{p \cdot (t - 1)}{1 - p}\right) 
\end{align} 
 
Pseudocode line \ref{line:map} implements this bound with conservative rounding towards positive infinity.  
When $\din > 0$ and no exception is raised in computing $\texttt{c} = \texttt{privacy\_map}(\din)$, then 
\[ 
    \ln\left(\frac{p \cdot (t - 1)}{1 - p}\right) \leq \texttt{c}. 
\] 
 
Therefore it has been shown that for every pair of elements $x, x' \in \texttt{input\_domain}$ and every $d_{DM}(x, x') \le \din$ with $\din \ge 0$,  
if $x, x'$ are $\din$-close then $\function(x),\function(x')$ are $\texttt{privacy\_map}(\din)$-close under $\texttt{output\_measure}$ (the Max-Divergence). 
\end{proof} 
 
\bibliographystyle{plain} 
\bibliography{randomized_response.bib} 
 
\end{document} 
